\documentclass[addpoints]{exam}
% \documentclass[addpoints,12pt]{exam}

\usepackage[slovene]{babel}
\usepackage{amsfonts}
\usepackage{amssymb}
\usepackage{amsmath}
\usepackage{float}
\usepackage{graphicx}
\usepackage{enumitem}
\usepackage{color}
\usepackage{eurosym}

\usepackage{pgfpages}
% \usepackage{tikz}
\usepackage{wrapfig}
\usepackage{pgfkeys}
\usepackage{pgfplots}
\usepackage{xcolor}
\usepackage{tkz-euclide}
\usepackage{xfp}




%Komentar naslednje vrstice glede na verzijo testa -- rešitve ali prostor za odgovore:
% \printanswers

% \linespread{2}


\pointsinrightmargin
\marginbonuspointname{}


\renewcommand{\solutiontitle}{\noindent\textbf{Rešitve in točkovanje:}\par\noindent}

\colorgrids
\definecolor{GridColor}{gray}{0.7}
\setlength{\gridsize}{7mm}

\extrawidth{5mm}
\setlength{\rightpointsmargin}{1.5cm}

\begin{document}

\runningfooter{}{}{}

\begin{center}
    \fbox{\fbox{\parbox{15cm}{\centering
    Popravljanje in pridobivanje ocen -- prvo pisno ocenjevanje znanja \\ 2. letnik -- test G \\ 23. januar 2026 \\ (Geometrija v ravnini)}}}
\end{center}

\vspace{0.1in}
\makebox[\textwidth]{Ime in priimek, oddelek:\enspace\hrulefill}


\begin{center}
    \hqword{Naloga}
    \hpword{Možne točke}
    % \chbpword{Dodatne točke}
    \hsword{Dosežene točke}
    \htword{Skupaj}  
    % \combinedpointtable[h][questions]
    \gradetable[h][questions]

\end{center}

\vspace{0.1in}


\begin{questions}

    % 1. vprašanje

    \question[5]
        Izračunajte velikosti vseh notranjih in zunanjih kotov trikotnika $\triangle ABC$, 
        če je vsota velikosti dveh zunanjih kotov $\alpha'+\beta'=234^\circ$, razlika velikosti dveh notranjih kotov pa $\alpha-\beta=28^\circ$.

    \begin{solution}[10cm]
        $ \dfrac{n(n-3)}{2}=44 \Rightarrow n^2-3n-88=0 \Rightarrow (n-11)(n+8)=0 \Rightarrow n=11 $ \\
        $ \varphi=\dfrac{n-2}{n}180^\circ=\dfrac{1620}{11}^\circ=147.\overline{27}^\circ $ \\ \\
        Za zapis in pravilno uporabo formule za število diagonal $1$ točka, za pravilen izračuna $1$ točka. 
        Za pravilen zapis in uporabo formule za izračun notranjega kota pravilnega $n$-kotnika $1$ točka, za pravilen izračun $1$ točka.
    \end{solution}


    % 2. vprašanje

    \question[3]
    Središčni kot je za $64^\circ$ večji od obodnega kota nad istim lokom. 
    Izračunajte velikosti obeh kotov.

    \begin{solution}[2cm]
        $\varphi+\psi=180^\circ$ \\
        $\varphi-\psi=25^\circ 29' 48''$ \\
        $\varphi=102^\circ 44' 54''$ in $\psi=77^\circ 15' 6 ''$ \\ \\
        Za pravilno postavitev enačb sistema po $1$ točka, reševanje sistema $1$ točka, za pravilen izračun rešitev po $1$ točka.
    \end{solution}

    

    \newpage
    % 3. vprašanje

    \question[6]
    Konstruirajte trapez s podatki $a=5~cm$, $e=f=4.5~cm$ in $\alpha=60^\circ$.

    \begin{solution}[15cm]
        Za zapis konstrukcije s pravilnim izrisom konstrukcije do $6$ točk. Zgolj za izris konstrukcije $1$ točka.
    \end{solution}




    \newpage
    % 4. vprašanje

    \question[5]
        S skice preberite ustrezne podatke ter izračunajte velikosti kotov $\alpha$, $\beta$, $\gamma$ in $\delta$.
        Pri tem velja, da so premice $p$, $q$ in $r$ vzporedne.
        \begin{figure}[H]
            \centering
            \begin{tikzpicture}
                % \clip (0,0) rectangle (14.000000,10.000000);
                {\footnotesize

                % Drawing line Q
                \draw [line width=0.016cm] (1.000000,3.000000) -- (5.800000,3.000000);%

                % Drawing line P
                \draw [line width=0.016cm] (1.000000,4.000000) -- (5.800000,4.000000);%

                % Drawing line R
                \draw [line width=0.016cm] (1.000000,1.500000) -- (5.800000,1.500000);%

                % Drawing line S
                \draw [line width=0.016cm] (2.426509,1.000000) -- (3.516142,4.800000);%

                % Drawing line T
                \draw [line width=0.016cm] (4.865030,1.000000) -- (1.321473,4.800000);%

                % Marking point p
                \draw (5.500000,4.000000) node [anchor=south] { $p$ };%

                % Marking point q
                \draw (5.500000,3.000000) node [anchor=south] { $q$ };%

                % Marking point r
                \draw (5.500000,1.500000) node [anchor=south] { $r$ };%

                % Marking point s
                \draw (3.500000,4.700000) node [anchor=west] { $s$ };%

                % Marking point t
                \draw (1.500000,4.700000) node [anchor=west] { $t$ };%

                % Marking point \alpha
                \draw (1.800000,4.000000) node [anchor=south east] { $\alpha$ };%

                % Marking point \beta
                \draw (3.300000,3.000000) node [anchor=north west] { $\beta$ };%

                % Marking point \gamma
                \draw (3.000000,3.000000) node [anchor=north east] { $\gamma$ };%

                % Marking point \delta
                \draw (4.450000,1.500000) node [anchor=south] { $\delta$ };%

                % Marking point {x-15^\circ}
                \draw (2.000000,3.000000) node [anchor=south] { ${x-15^\circ}$ };%

                % Marking point {2x-18^\circ}
                \draw (4.000000,4.000000) node [anchor=north] { ${2x-18^\circ}$ };%

                % Marking point {x+12^\circ}
                \draw (3.250000,1.500000) node [anchor=south] { ${x+12^\circ}$ };%
                }
            \end{tikzpicture}

        \end{figure}

    \begin{solution}[4cm]
        $x=65^\circ$ -- obodni kot nad lokom $AD$ \\
        $z=180^\circ-90^\circ-65^\circ=25^circ$ -- kot v pravokotnem trikotniku $\triangle ACD$ \\
        $y=180^\circ-65^\circ-25^\circ-35^\circ=55^\circ$ -- kot v trikotniku $\triangle ABD$ \\ \\
        Za izračun kota $x$ in utemeljitev $1$ točka, za izračun kota $y$ in utemeljitev $1$ točka, upoštevanje Talesovega izreka v polkrogu $1$ točka, izračuna kota $z$ $1$ točka.
    \end{solution}



    % 5. vprašanje

    \question[4]
        Izračunajte dolžine stranic pravokotnega trikotnika $\triangle ABC$, kjer sta $a$ in $b$ kateti, $c$ pa hipotenuza:
        $a=3x-1$, $b=5x$, $c=6x-1$. Določite dolžino pravokotne projekcije katete $a$ na hipotenuzo.


        \begin{solution}[9cm]
            Z upoštevanjem lastnosti pravokotnega trikotnika je $t_c=\frac{c}{2}=6~cm$. \\
            Za pravilen izračun dolžine $1$ točka.
        \end{solution}




    \newpage
    % 6. vprašanje

    \question[2]
        Marija stoji ponoči $5~m$ od vznožja $9~m$ visokega svetilnika. Njena senca je dolga $1~m$.
        Določite Marijino višino.


    \begin{solution}[11cm]
        Uporabimo lahko Pitagorov izrek za dolžini kraka in njegove pravokotne projekcije na osnovnico: $v^2=b^2-\left(\frac{a-c}{2}\right)^2=17^2-8^2=225 \Rightarrow v=15~cm$ \\ \\
        Za pravilen razmislek in izračun posameznih dolžin $2$ točki, pravilna uporaba Pitagorovega izreka $1$ točka, pravilen izračun $1$ točka.
    \end{solution}



    \question[3]
        Vsota števila stranic in diagonal $n$-kotnika je $105$. Izračunajte, kateri $n$-kotnik je to.

    % \newpage
    % dodatno vprašanje
    % \vspace{0.1in}
    % \makebox[\textwidth]{Ime in priimek, oddelek:\enspace\hrulefill}

    % ~

    % \bonusquestion[3]
    % \textit{Dodatna naloga:} Dokažite: Če iz točke, ki leži izven kroga, narišemo dve tangenti na krog, 
    % sta tangentna odseka (zveznici dane točke z dotikališčema) enako dolga.

    % \begin{solution}[1cm]
    %     Ob upoštevanju preloma droga dobimo pravokotni trikotnik s stranicami $7~m$, $20-x~m$ in $x~m$.
    %     Uporabimo pitagorov izrek $(20-x)^2=x^2+7^2 \Rightarrow 40x=351 \Rightarrow x=\frac{351}{40}~m$. \\ \\
    %     Pravilen razmislek in uporaba Pitagorovega izreka $1$ točka, reševanje enačbe $1$ točka, pravilen končni izračun $1$ točka.
    % \end{solution}


\end{questions}



\end{document}